\documentclass[a4paper, final]{article}

\usepackage[russian]{babel}
\usepackage{fontspec}
\setmainfont{Times New Roman}
\newfontfamily\cyrillicfonttt{Courier New}
\usepackage[14pt]{extsizes}
\usepackage[left=20mm, top=20mm, right=20mm, bottom=20mm, footskip=10mm]{geometry}
\usepackage{ragged2e}
\usepackage{indentfirst}
\usepackage{graphicx}
\usepackage{caption}
\usepackage{xcolor}
\usepackage{listings}
\usepackage{booktabs}
\usepackage{hyperref}
\usepackage{amsmath}
\usepackage{amssymb}
\usepackage{array}
\usepackage{multirow}

\definecolor{link_color}{HTML}{0000EE}
\hypersetup{colorlinks=true,linkcolor=link_color,urlcolor=link_color,citecolor=link_color}

\lstset{
    language=Python,
    basicstyle=\cyrillicfonttt\small,
    numbers=left,
    numberstyle=\tiny\color{gray},
    keywordstyle=\color{blue}\bfseries,
    commentstyle=\color{gray}\itshape,
    stringstyle=\color{red},
    frame=single,
    captionpos=b,
    breaklines=true,
    breakatwhitespace=false,
    showstringspaces=false,
    keepspaces=true,
    columns=flexible,
    fontadjust=true,
    basewidth={0.5em,0.45em},
    tabsize=4,
}

\begin{document}

{
\thispagestyle{empty}
\begin{center}
    \hfill\break\hfill\break
    МИНИСТЕРСТВО НАУКИ И ВЫСШЕГО ОБРАЗОВАНИЯ РОССИЙСКОЙ ФЕДЕРАЦИИ\\[5pt]
    Федеральное государственное автономное образовательное учреждение высшего образования
    \guillemotleft Санкт-Петербургский политехнический университет Петра Великого\guillemotright\\[10pt]
    Институт компьютерных наук и кибербезопасности\\[5pt]
    Направление: 02.03.01 Математика и компьютерные науки\\[30pt]
    {\large Отчет по дисциплине: \guillemotleft Вычислительная математика\guillemotright}\\[30pt]
    \LARGE\bf \guillemotleft Лабораторная работа 5\guillemotright.
\end{center}
\vspace{110pt}
\begin{tabular*}{460pt}{@{\extracolsep{\fill}} l r}
    Студент,\tabularnewline группы 5130201/40003 \hspace{50pt} & \line(1,0){100}\quad Селезнев К.\,Д.\\
    & \\
    Руководитель,\tabularnewline Преподаватель & \line(1,0){121.5}\quad Пак В.\,Г.\\
    & \\ & \\ & \\
    & \guillemotleft\line(1,0){30}\guillemotright\line(1,0){100}\,2026\,г.
\end{tabular*}\\
\vfill
\begin{center}Санкт-Петербург, 2026\end{center}
\newpage
}

\tableofcontents
\newpage

\section{Постановка задачи}

\textbf{Тема:} численное решение линейных систем.

\textbf{Цель:} изучить прямые и итерационные методы решения систем линейных уравнений.

\textbf{Задание:}
\begin{enumerate}
    \item Исследовать линейную систему методом Гаусса. Найти решение.
    \item Решить систему итерационными методами: Якоби и Зейделя. Вычислить пять итераций, оценить погрешность.
\end{enumerate}

Дана система $Ax = b$ (вариант~46):
\[
A = \begin{pmatrix}
4{,}03 & 0{,}21 & 0{,}12 & 0{,}09 \\
0{,}14 & 5{,}07 & 0{,}18 & 0{,}11 \\
0{,}08 & 0{,}15 & 6{,}04 & 0{,}13 \\
0{,}11 & 0{,}09 & 0{,}16 & 7{,}02
\end{pmatrix}, \qquad
b = \begin{pmatrix}
4{,}45 \\ 5{,}50 \\ 6{,}40 \\ 7{,}38
\end{pmatrix}.
\]

\newpage

\section{Задание 1. Метод Гаусса}

\subsection{Теоретические сведения}

Метод Гаусса~--- прямой метод решения системы $Ax = b$, состоящий из двух этапов.

\textbf{Прямой ход.} Расширенная матрица $[A\,|\,b]$ приводится к верхнетреугольному
виду. На шаге $k$ для каждой строки $i > k$ из строки $i$ вычитается строка $k$,
умноженная на множитель
\[
m_{ik} = \frac{a_{ik}^{(k)}}{a_{kk}^{(k)}}.
\]

\textbf{Обратный ход.} Из полученной треугольной системы находятся неизвестные:
\[
x_n = \frac{b_n^{(n)}}{a_{nn}^{(n)}}, \qquad
x_i = \frac{1}{a_{ii}^{(i)}}\!\left(b_i^{(i)} - \sum_{j=i+1}^{n} a_{ij}^{(i)}\,x_j\right),
\quad i = n-1, \ldots, 1.
\]

\subsection{Прямой ход}

Запишем расширенную матрицу $[A\,|\,b]$:
\[
\left(\begin{array}{cccc|c}
4{,}03 & 0{,}21 & 0{,}12 & 0{,}09 & 4{,}45 \\
0{,}14 & 5{,}07 & 0{,}18 & 0{,}11 & 5{,}50 \\
0{,}08 & 0{,}15 & 6{,}04 & 0{,}13 & 6{,}40 \\
0{,}11 & 0{,}09 & 0{,}16 & 7{,}02 & 7{,}38
\end{array}\right).
\]

\textbf{Шаг 1.} Ведущий элемент $a_{11} = 4{,}03$. Множители:
\[
m_{21} = \frac{0{,}14}{4{,}03} = 0{,}03474,\qquad
m_{31} = \frac{0{,}08}{4{,}03} = 0{,}01985,\qquad
m_{41} = \frac{0{,}11}{4{,}03} = 0{,}02730.
\]
Выполняем $R_i \leftarrow R_i - m_{i1} R_1$, $i = 2,3,4$:
\[
\left(\begin{array}{cccc|c}
4{,}03 & 0{,}21      & 0{,}12      & 0{,}09      & 4{,}45      \\
0      & 5{,}06270   & 0{,}17583   & 0{,}10687   & 5{,}34541   \\
0      & 0{,}14583   & 6{,}03762   & 0{,}12821   & 6{,}31166   \\
0      & 0{,}08427   & 0{,}15672   & 7{,}01754   & 7{,}25854
\end{array}\right).
\]

\textbf{Шаг 2.} Ведущий элемент $a_{22}^{(2)} = 5{,}06270$. Множители:
\[
m_{32} = \frac{0{,}14583}{5{,}06270} = 0{,}02881,\qquad
m_{42} = \frac{0{,}08427}{5{,}06270} = 0{,}01664.
\]
Выполняем $R_3 \leftarrow R_3 - m_{32} R_2$, $R_4 \leftarrow R_4 - m_{42} R_2$:
\[
\left(\begin{array}{cccc|c}
4{,}03 & 0{,}21    & 0{,}12      & 0{,}09      & 4{,}45      \\
0      & 5{,}06270 & 0{,}17583   & 0{,}10687   & 5{,}34541   \\
0      & 0         & 6{,}03255   & 0{,}12513   & 6{,}15769   \\
0      & 0         & 0{,}15380   & 7{,}01576   & 7{,}16956
\end{array}\right).
\]

\textbf{Шаг 3.} Ведущий элемент $a_{33}^{(3)} = 6{,}03255$. Множитель:
\[
m_{43} = \frac{0{,}15380}{6{,}03255} = 0{,}02549.
\]
Выполняем $R_4 \leftarrow R_4 - m_{43} R_3$:
\[
\left(\begin{array}{cccc|c}
4{,}03 & 0{,}21    & 0{,}12      & 0{,}09      & 4{,}45      \\
0      & 5{,}06270 & 0{,}17583   & 0{,}10687   & 5{,}34541   \\
0      & 0         & 6{,}03255   & 0{,}12513   & 6{,}15769   \\
0      & 0         & 0           & 7{,}01257   & 7{,}01257
\end{array}\right).
\]
Все диагональные элементы ненулевые~--- система имеет единственное решение.

\subsection{Обратный ход}

Из уравнения 4:
\[
x_4 = \frac{7{,}01257}{7{,}01257} = 1{,}00000.
\]

Из уравнения 3:
\[
x_3 = \frac{6{,}15769 - 0{,}12513 \cdot 1{,}00000}{6{,}03255}
     = \frac{6{,}15769 - 0{,}12513}{6{,}03255}
     = \frac{6{,}03256}{6{,}03255} = 1{,}00000.
\]

Из уравнения 2:
\[
x_2 = \frac{5{,}34541 - 0{,}17583 \cdot 1{,}00000 - 0{,}10687 \cdot 1{,}00000}{5{,}06270}
     = \frac{5{,}34541 - 0{,}17583 - 0{,}10687}{5{,}06270}= \]
     \[= \frac{5{,}06271}{5{,}06270} = 1{,}00000.
\]

Из уравнения 1:
\[
x_1 = \frac{4{,}45 - 0{,}21 \cdot 1{,}00000 - 0{,}12 \cdot 1{,}00000 - 0{,}09 \cdot 1{,}00000}{4{,}03} = \]
     \[= \frac{4{,}45 - 0{,}21 - 0{,}12 - 0{,}09}{4{,}03}
     = \frac{4{,}03}{4{,}03} = 1{,}00000.
\]

\textbf{Решение:}
\[
\boxed{x^* = (1{,}00000;\quad 1{,}00000;\quad 1{,}00000;\quad 1{,}00000).}
\]

\textbf{Проверка} $Ax^* = b$:
\begin{align*}
4{,}03\cdot1 + 0{,}21\cdot1 + 0{,}12\cdot1 + 0{,}09\cdot1 &= 4{,}45\ \checkmark\\
0{,}14\cdot1 + 5{,}07\cdot1 + 0{,}18\cdot1 + 0{,}11\cdot1 &= 5{,}50\ \checkmark\\
0{,}08\cdot1 + 0{,}15\cdot1 + 6{,}04\cdot1 + 0{,}13\cdot1 &= 6{,}40\ \checkmark\\
0{,}11\cdot1 + 0{,}09\cdot1 + 0{,}16\cdot1 + 7{,}02\cdot1 &= 7{,}38\ \checkmark
\end{align*}

\newpage

\section{Задание 2. Итерационные методы}

\subsection{Приведение системы к виду $x = Bx + c$}

Разрешим каждое уравнение системы $Ax = b$ относительно
соответствующего неизвестного:
\[
\left\{
\begin{aligned}
x_1 &= -\frac{a_{12}}{a_{11}}x_2 - \frac{a_{13}}{a_{11}}x_3 - \frac{a_{14}}{a_{11}}x_4
        + \frac{b_1}{a_{11}},\\
x_2 &= -\frac{a_{21}}{a_{22}}x_1 - \frac{a_{23}}{a_{22}}x_3 - \frac{a_{24}}{a_{22}}x_4
        + \frac{b_2}{a_{22}},\\
x_3 &= -\frac{a_{31}}{a_{33}}x_1 - \frac{a_{32}}{a_{33}}x_2 - \frac{a_{34}}{a_{33}}x_4
        + \frac{b_3}{a_{33}},\\
x_4 &= -\frac{a_{41}}{a_{44}}x_1 - \frac{a_{42}}{a_{44}}x_2
        - \frac{a_{43}}{a_{44}}x_3 + \frac{b_4}{a_{44}}.
\end{aligned}
\right.
\]

Коэффициенты $b_{ij} = -a_{ij}/a_{ii}$ (при $i \ne j$) и $c_i = b_i/a_{ii}$:
\begin{align*}
b_{12} &= -\tfrac{0{,}21}{4{,}03} = -0{,}05211, &
b_{13} &= -\tfrac{0{,}12}{4{,}03} = -0{,}02978, &
b_{14} &= -\tfrac{0{,}09}{4{,}03} = -0{,}02233, \\[4pt]
b_{21} &= -\tfrac{0{,}14}{5{,}07} = -0{,}02761, &
b_{23} &= -\tfrac{0{,}18}{5{,}07} = -0{,}03550, &
b_{24} &= -\tfrac{0{,}11}{5{,}07} = -0{,}02170, \\[4pt]
b_{31} &= -\tfrac{0{,}08}{6{,}04} = -0{,}01325, &
b_{32} &= -\tfrac{0{,}15}{6{,}04} = -0{,}02483, &
b_{34} &= -\tfrac{0{,}13}{6{,}04} = -0{,}02152, \\[4pt]
b_{41} &= -\tfrac{0{,}11}{7{,}02} = -0{,}01567, &
b_{42} &= -\tfrac{0{,}09}{7{,}02} = -0{,}01282, &
b_{43} &= -\tfrac{0{,}16}{7{,}02} = -0{,}02279.
\end{align*}

\[
c_1 = \tfrac{4{,}45}{4{,}03}  =  1{,}10422,\\[4pt]
c_2 = \tfrac{5{,}50}{5{,}07}  =  1{,}08481,\\[4pt]
c_3 = \tfrac{6{,}40}{6{,}04}  =  1{,}05960,\\[4pt]
c_4 = \tfrac{7{,}38}{7{,}02}  =  1{,}05128.
\]

Итерационная матрица $B$ и вектор $c$:
\[
B = \begin{pmatrix}
0          & -0{,}05211 & -0{,}02978 & -0{,}02233 \\
-0{,}02761 & 0          & -0{,}03550 & -0{,}02170 \\
-0{,}01325 & -0{,}02483 & 0          & -0{,}02152 \\
-0{,}01567 & -0{,}01282 & -0{,}02279 & 0
\end{pmatrix},\qquad
c = \begin{pmatrix}1{,}10422\\1{,}08481\\1{,}05960\\1{,}05128\end{pmatrix}.
\]

\subsection{Условие сходимости}

Норма матрицы $B$ (максимальная сумма по строке~--- норма $\|\cdot\|_\infty$):
\begin{align*}
\text{строка 1:}&\quad 0{,}05211+0{,}02978+0{,}02233 = 0{,}10422,\\
\text{строка 2:}&\quad 0{,}02761+0{,}03550+0{,}02170 = 0{,}08481,\\
\text{строка 3:}&\quad 0{,}01325+0{,}02483+0{,}02152 = 0{,}05960,\\
\text{строка 4:}&\quad 0{,}01567+0{,}01282+0{,}02279 = 0{,}05128.
\end{align*}
\[
\|B\|_\infty = \max\{0{,}10422;\;0{,}08481;\;0{,}05960;\;0{,}05128\} = 0{,}10422 < 1.
\]
По теореме~1 оба метода сходятся при любом начальном приближении.

\newpage

\subsection{Метод Якоби}

Расчётная формула метода Якоби~--- на каждой итерации используются
компоненты предыдущего приближения $x^{(k-1)}$:
\[
\left\{
\begin{aligned}
x_1^{(k)} &= b_{12}\,x_2^{(k-1)} + b_{13}\,x_3^{(k-1)} + b_{14}\,x_4^{(k-1)} + c_1,\\
x_2^{(k)} &= b_{21}\,x_1^{(k-1)} + b_{23}\,x_3^{(k-1)} + b_{24}\,x_4^{(k-1)} + c_2,\\
x_3^{(k)} &= b_{31}\,x_1^{(k-1)} + b_{32}\,x_2^{(k-1)} + b_{34}\,x_4^{(k-1)} + c_3,\\
x_4^{(k)} &= b_{41}\,x_1^{(k-1)} + b_{42}\,x_2^{(k-1)} + b_{43}\,x_3^{(k-1)} + c_4,
\end{aligned}
\right.\qquad k=1,2,\ldots
\]
Начальное приближение: $x^{(0)} = (0;\;0;\;0;\;0)^T$.

\medskip\textbf{Итерация $k=1$:}
\begin{align*}
x_1^{(1)} &= (-0{,}05211)\cdot0 + (-0{,}02978)\cdot0 + (-0{,}02233)\cdot0 + 1{,}10422 = \mathbf{1{,}10422},\\
x_2^{(1)} &= (-0{,}02761)\cdot0 + (-0{,}03550)\cdot0 + (-0{,}02170)\cdot0 + 1{,}08481 = \mathbf{1{,}08481},\\
x_3^{(1)} &= (-0{,}01325)\cdot0 + (-0{,}02483)\cdot0 + (-0{,}02152)\cdot0 + 1{,}05960 = \mathbf{1{,}05960},\\
x_4^{(1)} &= (-0{,}01567)\cdot0 + (-0{,}01282)\cdot0 + (-0{,}02279)\cdot0 + 1{,}05128 = \mathbf{1{,}05128}.
\end{align*}
$\|x^{(1)}-x^{(0)}\|_\infty = \max\{1{,}10422;\;1{,}08481;\;1{,}05960;\;1{,}05128\} = 1{,}10422.$

\medskip\textbf{Итерация $k=2$:}
\begin{align*}
x_1^{(2)} &= (-0{,}05211)\cdot1{,}08481 + (-0{,}02978)\cdot1{,}05960 + (-0{,}02233)\cdot1{,}05128 + 1{,}10422\\
           &= -0{,}05653 - 0{,}03155 - 0{,}02347 + 1{,}10422 = \mathbf{0{,}99266},\\
x_2^{(2)} &= (-0{,}02761)\cdot1{,}10422 + (-0{,}03550)\cdot1{,}05960 + (-0{,}02170)\cdot1{,}05128 + 1{,}08481\\
           &= -0{,}03049 - 0{,}03762 - 0{,}02281 + 1{,}08481 = \mathbf{0{,}99389},\\
x_3^{(2)} &= (-0{,}01325)\cdot1{,}10422 + (-0{,}02483)\cdot1{,}08481 + (-0{,}02152)\cdot1{,}05128 + 1{,}05960\\
           &= -0{,}01463 - 0{,}02694 - 0{,}02262 + 1{,}05960 = \mathbf{0{,}99541},\\
x_4^{(2)} &= (-0{,}01567)\cdot1{,}10422 + (-0{,}01282)\cdot1{,}08481 + (-0{,}02279)\cdot1{,}05960 + 1{,}05128\\
           &= -0{,}01730 - 0{,}01391 - 0{,}02415 + 1{,}05128 = \mathbf{0{,}99592}.
\end{align*}
$\|x^{(2)}-x^{(1)}\|_\infty = \max\{0{,}11156;\;0{,}09092;\;0{,}06419;\;0{,}05536\} = 0{,}11156.$

\medskip\textbf{Итерация $k=3$:}
\begin{align*}
x_1^{(3)} &= (-0{,}05211)\cdot0{,}99389 + (-0{,}02978)\cdot0{,}99541 + (-0{,}02233)\cdot0{,}99592 + 1{,}10422\\
           &= -0{,}05179 - 0{,}02964 - 0{,}02224 + 1{,}10422 = \mathbf{1{,}00055},\\
x_2^{(3)} &= (-0{,}02761)\cdot0{,}99266 + (-0{,}03550)\cdot0{,}99541 + (-0{,}02170)\cdot0{,}99592 + 1{,}08481\\
           &= -0{,}02741 - 0{,}03534 - 0{,}02161 + 1{,}08481 = \mathbf{1{,}00045},\\
x_3^{(3)} &= (-0{,}01325)\cdot0{,}99266 + (-0{,}02483)\cdot0{,}99389 + (-0{,}02152)\cdot0{,}99592 + 1{,}05960\\
           &= -0{,}01315 - 0{,}02468 - 0{,}02143 + 1{,}05960 = \mathbf{1{,}00034},\\
x_4^{(3)} &= (-0{,}01567)\cdot0{,}99266 + (-0{,}01282)\cdot0{,}99389 + (-0{,}02279)\cdot0{,}99541 + 1{,}05128\\
           &= -0{,}01556 - 0{,}01274 - 0{,}02269 + 1{,}05128 = \mathbf{1{,}00030}.
\end{align*}
$\|x^{(3)}-x^{(2)}\|_\infty = \max\{0{,}00789;\;0{,}00656;\;0{,}00493;\;0{,}00438\} = 0{,}00789.$

\medskip\textbf{Итерация $k=4$:}
\begin{align*}
x_1^{(4)} &= (-0{,}05211)\cdot1{,}00045 + (-0{,}02978)\cdot1{,}00034 + (-0{,}02233)\cdot1{,}00030 + 1{,}10422\\
           &= -0{,}05213 - 0{,}02979 - 0{,}02234 + 1{,}10422 = \mathbf{0{,}99996},\\
x_2^{(4)} &= (-0{,}02761)\cdot1{,}00055 + (-0{,}03550)\cdot1{,}00034 + (-0{,}02170)\cdot1{,}00030 + 1{,}08481\\
           &= -0{,}02763 - 0{,}03551 - 0{,}02171 + 1{,}08481 = \mathbf{0{,}99997},\\
x_3^{(4)} &= (-0{,}01325)\cdot1{,}00055 + (-0{,}02483)\cdot1{,}00045 + (-0{,}02152)\cdot1{,}00030 + 1{,}05960\\
           &= -0{,}01326 - 0{,}02484 - 0{,}02153 + 1{,}05960 = \mathbf{0{,}99998},\\
x_4^{(4)} &= (-0{,}01567)\cdot1{,}00055 + (-0{,}01282)\cdot1{,}00045 + (-0{,}02279)\cdot1{,}00034 + 1{,}05128\\
           &= -0{,}01568 - 0{,}01283 - 0{,}02280 + 1{,}05128 = \mathbf{0{,}99998}.
\end{align*}
$\|x^{(4)}-x^{(3)}\|_\infty = \max\{0{,}00059;\;0{,}00048;\;0{,}00036;\;0{,}00032\} = 0{,}00059.$

\medskip\textbf{Итерация $k=5$:}
\begin{align*}
x_1^{(5)} &= (-0{,}05211)\cdot0{,}99997 + (-0{,}02978)\cdot0{,}99998 + (-0{,}02233)\cdot0{,}99998 + 1{,}10422\\
           &= -0{,}05211 - 0{,}02978 - 0{,}02233 + 1{,}10422 = \mathbf{1{,}00000},\\
x_2^{(5)} &= (-0{,}02761)\cdot0{,}99996 + (-0{,}03550)\cdot0{,}99998 + (-0{,}02170)\cdot0{,}99998 + 1{,}08481\\
           &= -0{,}02761 - 0{,}03550 - 0{,}02170 + 1{,}08481 = \mathbf{1{,}00000},\\
x_3^{(5)} &= (-0{,}01325)\cdot0{,}99996 + (-0{,}02483)\cdot0{,}99997 + (-0{,}02152)\cdot0{,}99998 + 1{,}05960\\
           &= -0{,}01325 - 0{,}02483 - 0{,}02152 + 1{,}05960 = \mathbf{1{,}00000},\\
x_4^{(5)} &= (-0{,}01567)\cdot0{,}99996 + (-0{,}01282)\cdot0{,}99997 + (-0{,}02279)\cdot0{,}99998 + 1{,}05128\\
           &= -0{,}01567 - 0{,}01282 - 0{,}02279 + 1{,}05128 = \mathbf{1{,}00000}.
\end{align*}
$\|x^{(5)}-x^{(4)}\|_\infty = \max\{0{,}00004;\;0{,}00003;\;0{,}00002;\;0{,}00002\} = 0{,}00004.$

\begin{table}[h!]
\centering
\caption{Итерации метода Якоби}
\label{tab:jacobi}
\renewcommand{\arraystretch}{1.3}
\begin{tabular}{|c|c|c|c|c|c|c|}
\hline
 & $x^{(0)}$ & $x^{(1)}$ & $x^{(2)}$ & $x^{(3)}$ & $x^{(4)}$ & $x^{(5)}$ \\
\hline
$x_1$ & 0 & 1,10422 & 0,99266 & 1,00055 & 0,99996 & 1,00000 \\
$x_2$ & 0 & 1,08481 & 0,99389 & 1,00045 & 0,99997 & 1,00000 \\
$x_3$ & 0 & 1,05960 & 0,99541 & 1,00034 & 0,99998 & 1,00000 \\
$x_4$ & 0 & 1,05128 & 0,99592 & 1,00030 & 0,99998 & 1,00000 \\
\hline
$\|x^{(k)}-x^{(k-1)}\|_\infty$ & --- & 1,10422 & 0,11156 & 0,00789 & 0,00059 & 0,00004 \\
\hline
\end{tabular}
\end{table}

Апостериорная оценка погрешности (теорема~2):
\[
\Delta x^{(5)} \leq \frac{\|B\|}{1-\|B\|}\cdot\|x^{(5)}-x^{(4)}\|_\infty
= \frac{0{,}10422}{1-0{,}10422}\cdot 0{,}00004
= 0{,}11635\cdot 0{,}00004 = 0{,}000005.
\]
Фактическая погрешность: $\|x^{(5)}-x^*\|_\infty \approx 3\cdot10^{-6}.$

\newpage

\subsection{Метод Зейделя}

В методе Зейделя при вычислении $x_i^{(k)}$ используются
уже найденные на шаге $k$ значения $x_1^{(k)},\ldots,x_{i-1}^{(k)}$.

Матрица $B$ разбивается: $B = B_1 + B_2$, где
\[
B_1 = \begin{pmatrix}
0          & 0          & 0          & 0 \\
-0{,}02761 & 0          & 0          & 0 \\
-0{,}01325 & -0{,}02483 & 0          & 0 \\
-0{,}01567 & -0{,}01282 & -0{,}02279 & 0
\end{pmatrix} \]
\[
B_2 = \begin{pmatrix}
0 & -0{,}05211 & -0{,}02978 & -0{,}02233 \\
0 & 0          & -0{,}03550 & -0{,}02170 \\
0 & 0          & 0          & -0{,}02152 \\
0 & 0          & 0          & 0
\end{pmatrix}.
\]
\[
\|B_1\|_\infty = 0{,}05128,\qquad \|B_2\|_\infty = 0{,}10422,\qquad
\|B_1\|_\infty+\|B_2\|_\infty = 0{,}15550 < 1.
\]
По теореме~3 метод сходится. Скорость сходимости:
\[
q = \frac{\|B_2\|}{1-\|B_1\|} = \frac{0{,}10422}{1-0{,}05128} = 0{,}10985.
\]

Расчётная формула метода Зейделя:
\[
\left\{
\begin{aligned}
x_1^{(k)} &= b_{12}\,x_2^{(k-1)} + b_{13}\,x_3^{(k-1)} + b_{14}\,x_4^{(k-1)} + c_1,\\
x_2^{(k)} &= b_{21}\,x_1^{(k)} + b_{23}\,x_3^{(k-1)} + b_{24}\,x_4^{(k-1)} + c_2,\\
x_3^{(k)} &= b_{31}\,x_1^{(k)} + b_{32}\,x_2^{(k)} + b_{34}\,x_4^{(k-1)} + c_3,\\
x_4^{(k)} &= b_{41}\,x_1^{(k)} + b_{42}\,x_2^{(k)} + b_{43}\,x_3^{(k)} + c_4.
\end{aligned}
\right.\qquad k=1,2,\ldots
\]
Начальное приближение: $x^{(0)} = (0;\;0;\;0;\;0)^T$.

\medskip\textbf{Итерация $k=1$:}
\begin{align*}
x_1^{(1)} &= (-0{,}05211)\cdot0 + (-0{,}02978)\cdot0 + (-0{,}02233)\cdot0 + 1{,}10422 = \mathbf{1{,}10422},\\
x_2^{(1)} &= (-0{,}02761)\cdot\mathbf{1{,}10422} + (-0{,}03550)\cdot0 + (-0{,}02170)\cdot0 + 1{,}08481\\
           &= -0{,}03049 + 1{,}08481 = \mathbf{1{,}05432},\\
x_3^{(1)} &= (-0{,}01325)\cdot\mathbf{1{,}10422} + (-0{,}02483)\cdot\mathbf{1{,}05432} + (-0{,}02152)\cdot0 + 1{,}05960\\
           &= -0{,}01463 - 0{,}02618 + 1{,}05960 = \mathbf{1{,}01879},\\
x_4^{(1)} &= (-0{,}01567)\cdot\mathbf{1{,}10422} + (-0{,}01282)\cdot\mathbf{1{,}05432}
             + (-0{,}02279)\cdot\mathbf{1{,}01879} + 1{,}05128\\
           &= -0{,}01730 - 0{,}01352 - 0{,}02322 + 1{,}05128 = \mathbf{0{,}99724}.
\end{align*}
$\|x^{(1)}-x^{(0)}\|_\infty = \max\{1{,}10422;\;1{,}05432;\;1{,}01879;\;0{,}99724\} = 1{,}10422.$

\medskip\textbf{Итерация $k=2$:}
\begin{align*}
x_1^{(2)} &= (-0{,}05211)\cdot1{,}05432 + (-0{,}02978)\cdot1{,}01879 + (-0{,}02233)\cdot0{,}99724 + 1{,}10422\\
           &= -0{,}05494 - 0{,}03034 - 0{,}02227 + 1{,}10422 = \mathbf{0{,}99667},\\
x_2^{(2)} &= (-0{,}02761)\cdot\mathbf{0{,}99667} + (-0{,}03550)\cdot1{,}01879 + (-0{,}02170)\cdot0{,}99724 + 1{,}08481\\
           &= -0{,}02752 - 0{,}03617 - 0{,}02164 + 1{,}08481 = \mathbf{0{,}99949},\\
x_3^{(2)} &= (-0{,}01325)\cdot\mathbf{0{,}99667} + (-0{,}02483)\cdot\mathbf{0{,}99949} + (-0{,}02152)\cdot0{,}99724 + 1{,}05960\\
           &= -0{,}01320 - 0{,}02482 - 0{,}02146 + 1{,}05960 = \mathbf{1{,}00012},\\
x_4^{(2)} &= (-0{,}01567)\cdot\mathbf{0{,}99667} + (-0{,}01282)\cdot\mathbf{0{,}99949}
             + (-0{,}02279)\cdot\mathbf{1{,}00012} + 1{,}05128\\
           &= -0{,}01562 - 0{,}01281 - 0{,}02279 + 1{,}05128 = \mathbf{1{,}00006}.
\end{align*}
$\|x^{(2)}-x^{(1)}\|_\infty = \max\{0{,}10755;\;0{,}05483;\;0{,}01867;\;0{,}00282\} = 0{,}10755.$

\medskip\textbf{Итерация $k=3$:}
\begin{align*}
x_1^{(3)} &= (-0{,}05211)\cdot0{,}99949 + (-0{,}02978)\cdot1{,}00012 + (-0{,}02233)\cdot1{,}00006 + 1{,}10422\\
           &= -0{,}05208 - 0{,}02978 - 0{,}02233 + 1{,}10422 = \mathbf{1{,}00002},\\
x_2^{(3)} &= (-0{,}02761)\cdot\mathbf{1{,}00002} + (-0{,}03550)\cdot1{,}00012 + (-0{,}02170)\cdot1{,}00006 + 1{,}08481\\
           &= -0{,}02761 - 0{,}03550 - 0{,}02170 + 1{,}08481 = \mathbf{0{,}99999},\\
x_3^{(3)} &= (-0{,}01325)\cdot\mathbf{1{,}00002} + (-0{,}02483)\cdot\mathbf{0{,}99999} + (-0{,}02152)\cdot1{,}00006 + 1{,}05960\\
           &= -0{,}01325 - 0{,}02483 - 0{,}02152 + 1{,}05960 = \mathbf{1{,}00000},\\
x_4^{(3)} &= (-0{,}01567)\cdot\mathbf{1{,}00002} + (-0{,}01282)\cdot\mathbf{0{,}99999}
             + (-0{,}02279)\cdot\mathbf{1{,}00000} + 1{,}05128\\
           &= -0{,}01567 - 0{,}01282 - 0{,}02279 + 1{,}05128 = \mathbf{1{,}00000}.
\end{align*}
$\|x^{(3)}-x^{(2)}\|_\infty = \max\{0{,}00335;\;0{,}00050;\;0{,}00012;\;0{,}00006\} = 0{,}00335.$

\medskip\textbf{Итерации $k=4,5$} практически совпадают с точным решением
(изменения $< 10^{-5}$):
\[
x^{(4)} \approx x^{(5)} \approx (1{,}00000;\;1{,}00000;\;1{,}00000;\;1{,}00000)^T.
\]

\begin{table}[h!]
\centering
\caption{Итерации метода Зейделя}
\label{tab:seidel}
\renewcommand{\arraystretch}{1.3}
\begin{tabular}{|c|c|c|c|c|c|c|}
\hline
 & $x^{(0)}$ & $x^{(1)}$ & $x^{(2)}$ & $x^{(3)}$ & $x^{(4)}$ & $x^{(5)}$ \\
\hline
$x_1$ & 0 & 1,10422 & 0,99667 & 1,00002 & 1,00000 & 1,00000 \\
$x_2$ & 0 & 1,05432 & 0,99949 & 0,99999 & 1,00000 & 1,00000 \\
$x_3$ & 0 & 1,01879 & 1,00012 & 1,00000 & 1,00000 & 1,00000 \\
$x_4$ & 0 & 0,99724 & 1,00006 & 1,00000 & 1,00000 & 1,00000 \\
\hline
$\|x^{(k)}-x^{(k-1)}\|_\infty$ & --- & 1,10422 & 0,10755 & 0,00335 & $<10^{-4}$ & $<10^{-5}$ \\
\hline
\end{tabular}
\end{table}

Апостериорная оценка погрешности (теорема~4):
\[
\Delta x^{(3)} \leq \frac{\|B_2\|}{1-\|B\|}\cdot\|x^{(3)}-x^{(2)}\|_\infty
= \frac{0{,}10422}{1-0{,}10422}\cdot 0{,}00335
= 0{,}11635\cdot 0{,}00335 = 0{,}000390.
\]
Фактическая погрешность: $\|x^{(3)}-x^*\|_\infty \approx 2\cdot10^{-5}.$

\newpage

\section{Сравнение методов}

\begin{table}[h!]
\centering
\caption{Сравнение итерационных методов}
\label{tab:compare}
\renewcommand{\arraystretch}{1.4}
\begin{tabular}{|l|c|c|}
\hline
Показатель & Якоби & Зейделя \\
\hline
$\|B\|_\infty$ & $0{,}10422$ & $0{,}10422$ \\
\hline
$\|B_1\|_\infty + \|B_2\|_\infty$ & --- & $0{,}15550 < 1$ \\
\hline
$\|x^{(3)}-x^*\|_\infty$ & $5\cdot10^{-4}$ & $2\cdot10^{-5}$ \\
\hline
$\|x^{(5)}-x^*\|_\infty$ & $3\cdot10^{-6}$ & $<10^{-6}$ \\
\hline
\end{tabular}
\end{table}

Оба метода сходятся очень быстро благодаря сильному диагональному преобладанию
матрицы $A$: норма $\|B\|_\infty = 0{,}104 \ll 1$. Метод Зейделя сходится
быстрее: после 3~итераций его погрешность на порядок меньше, чем у Якоби.
Это объясняется тем, что при вычислении каждой компоненты $x_i^{(k)}$
немедленно используются уточнённые значения предыдущих компонент.

\newpage

\section*{Приложение}
\addcontentsline{toc}{section}{Приложение}

\begin{lstlisting}[caption={Метод Гаусса и итерационные методы}, label=lst:code]
import numpy as np

A = np.array([
    [4.03, 0.21, 0.12, 0.09],
    [0.14, 5.07, 0.18, 0.11],
    [0.08, 0.15, 6.04, 0.13],
    [0.11, 0.09, 0.16, 7.02]
])
b = np.array([4.45, 5.50, 6.40, 7.38])
n = 4

# ===== Metod Gaussa =====
Ab = np.hstack([A.astype(float), b.reshape(-1,1)])
for k in range(n):
    for i in range(k+1, n):
        m = Ab[i,k] / Ab[k,k]
        Ab[i] -= m * Ab[k]
x = np.zeros(n)
for i in range(n-1, -1, -1):
    x[i] = (Ab[i,n] - Ab[i,i+1:n] @ x[i+1:n]) / Ab[i,i]
print("Gauss:", np.round(x, 5))

# ===== Matritsa B i vektor c =====
B = np.zeros((n,n))
c = np.zeros(n)
for i in range(n):
    c[i] = b[i] / A[i,i]
    for j in range(n):
        if i != j:
            B[i,j] = -A[i,j] / A[i,i]
print("||B||_inf =", np.max(np.sum(np.abs(B), axis=1)))

# ===== Metod Yakobi =====
x_j = np.zeros(n)
for k in range(1, 6):
    x_j_new = B @ x_j + c
    diff = np.max(np.abs(x_j_new - x_j))
    print(f"Jacobi k={k}: {np.round(x_j_new,5)}, diff={diff:.5f}")
    x_j = x_j_new

# ===== Metod Zeydelya =====
x_s = np.zeros(n)
for k in range(1, 6):
    x_s_new = x_s.copy()
    for i in range(n):
        x_s_new[i] = c[i]
        for j in range(i):
            x_s_new[i] += B[i,j] * x_s_new[j]  # novye znacheniya
        for j in range(i+1, n):
            x_s_new[i] += B[i,j] * x_s[j]       # starye znacheniya
    diff = np.max(np.abs(x_s_new - x_s))
    print(f"Seidel k={k}: {np.round(x_s_new,5)}, diff={diff:.5f}")
    x_s = x_s_new
\end{lstlisting}

\end{document}